\documentclass[11pt,a4paper]{article}
\usepackage[T1]{fontenc}
\usepackage{lmodern}
\usepackage[margin=24mm]{geometry}
\usepackage{amsmath,amssymb,amsthm}
\usepackage{microtype}
\usepackage{xcolor}
\usepackage{listings}
\usepackage{fancyhdr}
\usepackage[hidelinks]{hyperref}
\definecolor{ink}{HTML}{17334A}
\definecolor{shade}{HTML}{F4F6F8}
\newtheorem{proposition}{Proposition}
\newtheorem{lemma}[proposition]{Lemma}
\newtheorem{corollary}[proposition]{Corollary}
\DeclareMathOperator{\rad}{rad}
\pagestyle{fancy}
\fancyhf{}
\fancyhead[L]{\small\sffamily Integer solutions of $x^4y^3=z^2+1$}
\fancyhead[R]{\small\sffamily Mathematical analysis}
\fancyfoot[C]{\small\thepage}
\setlength{\headheight}{15pt}
\setlength{\parindent}{0pt}
\setlength{\parskip}{5pt}
\setlength{\emergencystretch}{2em}
\lstset{language=Python,basicstyle=\ttfamily\footnotesize,
  backgroundcolor=\color{shade},frame=single,rulecolor=\color{black!15},
  columns=fullflexible,keepspaces=true,showstringspaces=false,
  xleftmargin=5pt,xrightmargin=5pt,aboveskip=8pt,belowskip=8pt}
\hypersetup{pdftitle={On a Diophantine equation with fourth, third, and second powers},
  pdfsubject={Self-contained restrictions, a polynomial obstruction, and conditional finiteness}}

\begin{document}
\thispagestyle{empty}
{\Large\bfseries\color{ink} On the equation $x^4y^3=z^2+1$}\par
{\large Elementary restrictions and the obstacle to infinitude}\par
\vspace{3mm}

\textbf{Status of the requested conclusion.}
I have not established infinitely many integer solutions. The arguments below
prove several unconditional restrictions, including the impossibility of a
nonconstant polynomial family. They also prove that the $abc$ conjecture would
imply \emph{finitely} many integer solutions. Consequently, an infinitude proof
would disprove that conjecture. This is not an unconditional disproof of
infinitude. Nor is a complete classification established here.

The solutions $(1,1,0)$ and $(-1,1,0)$ are immediate. All assertions below are
proved from the stated assumptions; no internet search or external mathematical
sources were used.

\section*{1. Intuition: rational formulas versus integer solutions}
An apparently promising substitution is $u=xy$. It gives the exact rational
parametrization
\[
 x=\frac{t^2+1}{u^3},\qquad
 y=\frac{u^4}{t^2+1},\qquad z=t,
 \qquad (u,t\in\mathbb Q,\ u\ne0).
\]
Indeed, multiplying the fourth power of the first expression by the cube of
the second gives $t^2+1$. Conversely, every rational solution is obtained by
taking $u=xy$ and $t=z$.

For \emph{integers} $u,t$, however, integrality of both coordinates requires
\[
 u^3\mid t^2+1\qquad\text{and}\qquad t^2+1\mid u^4.
\]
Thus the rational parametrization leaves the essential arithmetic difficulty
intact. Equivalently, in any integer solution, each prime exponent in $z^2+1$
has the form $4a+3b$ with $a,b\ge0$. A positive exponent is therefore at least
$3$; exponents $1,2,5$ are impossible. This scarcity of distinct prime factors
is the reason the $abc$ obstruction in Section~4 appears.

\begin{proposition}[Elementary restrictions]
Every integer solution has $x\ne0$, $y>0$, $x,y$ odd, and $z$ even. Every prime
dividing $xy$ is $1\pmod4$, and $y\equiv1$ or $13\pmod{16}$.
If $y$ is a square, the solution is $(\pm1,1,0)$.
\end{proposition}
\begin{proof}
Positivity of $z^2+1$ forces $x\ne0$ and $y>0$. An even $x$ or an even $y$
would give $z^2\equiv-1\pmod4$, which is impossible. Therefore $x,y$ are odd
and $z$ is even. If an odd prime $p\mid xy$, then $z^2\equiv-1\pmod p$.
Multiplication by $z$ on the nonzero residues modulo $p$ partitions them into
orbits of length exactly $4$: its square is multiplication by $-1$, and its
fourth power is the identity. Hence $4\mid p-1$.
Modulo $16$, $x^4\equiv1$ and $z^2+1\equiv1$ or $5$.
Cubing permutes the odd residues and is its own inverse, so $y\equiv1$ or $13$.
Finally, if $y=v^2$ with $v>0$, put $w=x^2v^3>0$. Then
$(w-z)(w+z)=1$, forcing $w=1$, $z=0$, and $|x|=v=1$.
\end{proof}

\newpage
\section*{2. Another infinite class of candidates is excluded}
The following elementary cube obstruction also rules out every $|x|$ that is
a perfect cube. We include its proof so this restriction is self-contained.

\begin{lemma}
The integer equation $t^2+1=v^3$ has only the solution $(t,v)=(0,1)$.
\end{lemma}
\begin{proof}
Necessarily $v>0$. If $t$ were odd, $t^2+1\equiv2\pmod4$, whereas a cube is
$0,1$, or $3\pmod4$. Hence $t$ is even and $v$ is odd.

We use unique factorization in $\mathbb Z[i]$, justified as follows. Its norm
is $N(a+bi)=a^2+b^2$. Rounding both coordinates of $\alpha/\beta$ to nearest
integers gives a Gaussian integer $q$ with
$N(\alpha-q\beta)\le N(\beta)/2<N(\beta)$. Thus the Euclidean algorithm
and B\'ezout's identity hold. B\'ezout's identity makes every irreducible
prime; induction on the norm gives factorization, and primality gives
uniqueness.

In $(t+i)(t-i)=v^3$, the two factors are coprime: a common divisor divides
$2i$, so its norm divides $4$, while it also divides the odd integer $t^2+1$.
Its norm is therefore $1$. Unique factorization now shows that
$t+i=\varepsilon(a+bi)^3$ for a unit $\varepsilon$. Cubing permutes the four
units $1,-1,i,-i$, so the unit can be absorbed into the cube. Comparing
imaginary parts gives
\[
 b(3a^2-b^2)=1.
\]
Thus $b=\pm1$. The case $b=1$ requires $3a^2=2$, which is impossible;
the case $b=-1$ gives $a=0$. Consequently $t=0$ and $v=1$.
\end{proof}

\begin{corollary}
If $|x|$ is a perfect cube, then $(x,y,z)=(\pm1,1,0)$.
\end{corollary}
\begin{proof}
Write $|x|=a^3$ with $a>0$. Then $z^2+1=(a^4y)^3$. The lemma yields
$z=0$ and $a^4y=1$, so $a=y=1$.
\end{proof}

\section*{3. Why a polynomial construction cannot work}
\begin{proposition}
If $X,Y,Z\in\mathbb Q[T]$ satisfy $X^4Y^3=Z^2+1$ identically, all three
polynomials are constant.
\end{proposition}
\begin{proof}
$X$ and $Y$ are nonzero. If $Z$ is constant, comparison of degrees proves
the claim. Otherwise write $a=\deg X$, $b=\deg Y$, $c=\deg Z\ge1$.
Comparison of degrees gives $2c=4a+3b$. Differentiation gives
\[
 2ZZ'=X^3Y^2(4X'Y+3XY').
\]
No nonconstant polynomial divides both $Z$ and $XY$: such a divisor would
divide $Z^2+1$ and $Z$, hence $1$. Therefore $X^3Y^2$ is coprime to $Z$,
and the displayed identity implies $X^3Y^2\mid Z'$. Taking degrees yields
\[
 3a+2b\le c-1=2a+\tfrac32b-1,
 \qquad\text{so}\qquad a+\tfrac12b\le-1,
\]
contradicting $a,b\ge0$.
\end{proof}

The intuition is that differentiation preserves too much of the repeated
factors on the left, while reducing the available degree on the right.

\newpage
\begin{corollary}[Rational formulas at integer parameters]
No nonconstant triple of rational functions in $T$ can yield integer solutions
for infinitely many integer values of $T$.
\end{corollary}
\begin{proof}
First suppose $f\in\mathbb Q(T)$ is integer-valued at infinitely many integers.
Polynomial division writes $f=P+R/Q$, with $P,R,Q\in\mathbb Q[T]$,
$Q\ne0$, and either $R=0$ or $\deg R<\deg Q$. Choose a positive integer $D$
such that $DP\in\mathbb Z[T]$. Along those integer inputs, which are unbounded
in absolute value, $Df-DP=DR/Q$ is an integer tending to zero. It therefore
vanishes at infinitely many inputs, forcing $R=0$. Thus $f$ is a polynomial.

Apply this argument to each coordinate. The equation holds identically,
because a nonzero rational function has only finitely many zeros. The preceding
proposition then makes every coordinate constant.
\end{proof}

This result concerns rational functions evaluated at \emph{integer} parameters.
It does not rule out recurrence sequences or other possible constructions of
integer solutions, and it does not conflict with the two-parameter rational
parametrization in Section~1.

\section*{4. Infinitude would contradict the $abc$ conjecture}
For a positive integer $m$, let $\rad(m)$ be the product of its distinct prime
divisors, with $\rad(1)=1$. The form of the \emph{$abc$ conjecture} needed here
says that for every $\varepsilon>0$ there is a constant $K_\varepsilon$ such that
pairwise coprime positive integers $A,B,C$ with $A+B=C$ satisfy
\[
 C\le K_\varepsilon\rad(ABC)^{1+\varepsilon}.
\]
This is an additional conjectural hypothesis, not an ingredient in any of the
preceding unconditional results.

\begin{proposition}[Conditional finiteness]
If the $abc$ conjecture holds, $x^4y^3=z^2+1$ has only finitely many integer
solutions.
\end{proposition}
\begin{proof}
The case $z=0$ gives only $(\pm1,1,0)$. For $z\ne0$, put
$C=z^2+1=x^4y^3$. Every prime exponent in $C$ is at least $3$, so
\[
 \rad(C)^3\le C.
\]
Since $|z|<C^{1/2}$, we have
\[
 \rad(z^2C)\le |z|\rad(C)\le C^{1/2+1/3}=C^{5/6}.
\]
Apply the conjecture to $1+z^2=C$, with $\varepsilon=1/10$. These three
integers are pairwise coprime, and thus
\[
 C\le K_{1/10}\rad(z^2C)^{11/10}
   \le K_{1/10}C^{11/12},
 \qquad C\le K_{1/10}^{12}.
\]
This bounds $z$. Each fixed positive $C$ has only finitely many factorizations
$C=x^4y^3$ with $x\ne0$, $y>0$, proving the assertion.
\end{proof}

The argument explains the obstruction quantitatively: $z^2+1$ would be large
relative to the product of the distinct primes in $z^2(z^2+1)$. It supplies no
unconditional global bound and no usable numerical bound without control of
$K_{1/10}$.

\newpage
\section*{5. A reproducible finite check}
An exact integer computation gives the following bounded result:
\[
 \boxed{\ |z|\le10^{10}\quad\Longrightarrow\quad
 (x,y,z)\in\{(1,1,0),(-1,1,0)\}.\ }
\]
This is a computational verification within the stated bound, not a claim
about larger solutions. The following complete Python code uses only integer
arithmetic and the standard library; it checked $2{,}503{,}682$ candidate pairs.

\begin{lstlisting}
from math import isqrt

def icbrt(n):
    # The greatest integer r such that r**3 <= n.
    lo, hi = 0, 1 << ((n.bit_length() + 2) // 3)
    while lo + 1 < hi:
        mid = (lo + hi) // 2
        if mid**3 <= n:
            lo = mid
        else:
            hi = mid
    return lo

B = 10**10
C = B*B + 1
solutions = []
checked = 0
for x in range(1, isqrt(isqrt(C)) + 1, 2):
    ymax = icbrt(C // x**4)
    for y in range(1, ymax + 1, 4):
        n = x**4 * y**3 - 1
        z = isqrt(n)
        checked += 1
        if z*z == n:
            solutions.append((x, y, z))
print(checked, solutions)
# Output: 2503682 [(1, 1, 0)]
\end{lstlisting}

\textbf{Why the check is exhaustive within its bound.}
Sign changes of $x$ or $z$ preserve the equation, so it suffices to take $x>0$
and $z\ge0$. If $z\le B$, then $x^4y^3\le B^2+1=C$. Since $y\ge1$,
\[
 1\le x\le\lfloor C^{1/4}\rfloor,
 \qquad
 1\le y\le\left\lfloor\sqrt[3]{\lfloor C/x^4\rfloor}\right\rfloor.
\]
The outer loop uses $\lfloor C^{1/4}\rfloor=\operatorname{isqrt}
(\operatorname{isqrt}(C))$. The function \texttt{icbrt} maintains
$\texttt{lo}^3\le n<\texttt{hi}^3$, terminates when the endpoints are adjacent,
and therefore returns the exact integer cube root. The parity and congruence
restrictions justify testing only odd $x$ and $y\equiv1\pmod4$. Finally,
\texttt{isqrt} and the test $z^2=n$ detect squares exactly. Thus every possible
solution in the range is tested.

\textbf{What remains unproved.}
This note establishes neither global finiteness nor global infinitude, and it
does not prove that $(\pm1,1,0)$ are the complete list without a bound on $z$.
The requested unconditional proof of infinitely many integer solutions is
therefore not supplied. The conditional obstruction and the unconditional
restrictions above should not be mistaken for such a proof.

\end{document}
